\documentclass[11pt]{exam}


\usepackage[margin=1in, nohead, pdftex]{geometry}
\usepackage{multicol}

\def\R{\mathbb{R}}

\topmargin -0.2in

%\usepackage[pdftex, colorlinks=true, linkcolor=blue, citecolor=blue, urlcolor=blue, 
%			letterpaper, pdftitle={}, pdfpagemode=None, pdfauthor={Justin Sukiennik}]{hyperref}

\parindent 0ex

%=============================================================================================================
\begin{document}
\pagestyle{empty}
%=============================================================================================================
\textbf{MATH 96} \hspace{\stretch{1}} \textbf{Name:} \makebox[2.77in]{\hrulefill}\\
\textbf{Fall 2024} \\
\textbf{Quiz 1 Version A}\\
\textbf{20 Minutes}  \hspace{\stretch{1}} 

\rule[1ex]{\textwidth}{.1pt}
You may \emph{not} use textbooks, course notes, or a calculator during this exam. Organize your work in a reasonable, neat, and coherent way.  Any answer must be supported by calculations, explanation, and/or algebraic work to receive full credit. Partial credit may be given.
\vspace{5pt}\\
\rule[1ex]{\textwidth}{.1pt}
%=============================================================================================================
\noprintanswers
\begin{questions}
%=======================================================================================================================
\vspace{-10pt}
\question[8] Simplify each expression fully.
\noaddpoints
\begin{multicols}{2}
\begin{parts}
\part $\displaystyle 3(-4n + 2) + 6(n - 1)$
\vspace{2.75in}
\part $\displaystyle \frac{(-3)^3-2^3}{-2^3 + 3^2 +4}$
\end{parts}
\addpoints
\end{multicols}
\vspace{2.75in}
\question[22] Solve each linear equation. \emph{Show your work for credit.}
\noaddpoints
\begin{parts}
\part[6] $ -2x +3 = 9$
\newpage
%=============================================================================================================
\part[8] $3 (2 - 4x) = 2x- 8$
\vspace{4.25in}
\part[8] $\displaystyle \frac{1}{4} + \left( -\frac{2}{3} x  \right)= \frac{3}{12}$
\end{parts}
\addpoints
%=============================================================================================================
\end{questions}
\end{document}
